\section{导读：机器人创业为什么“不浪漫”}

本节先把阅读方式定下来：不要把它当成“谁离开了哪家公司”的新闻八卦，也不要只当成星海图的公司宣传。更有价值的读法，是把它看作一个创业者如何把自动驾驶训练迁移到具身智能的口述案例。

这期访谈表面上是星海图创始人高继扬的个人经历，实际是一堂关于“物理世界 AI 公司怎么长出来”的产业课。主持人一开始提出一个疑惑：为什么中国具身智能行业里，很少出现像大模型创业那样带有浓厚技术浪漫主义的人？高继扬给出的答案贯穿全场：机器人链条极长，周期极长，团队天然要把头伸到土里。

这里的“土里”不是消极词，而是物理世界的真实约束：整机、供应链、数据采集、客户现场、端侧延迟、维修、渠道、融资和组织成长。大模型或软件应用可以先从模型和传播启动，机器人公司却很难绕开硬件和客户价值。因此，这份笔记把本期整理成“物理世界 AI 的产业训练营”：Waymo 教他系统工程，Momenta 教他量产交付，星海图把整机、数据、模型和客户闭环重新组合。

\begin{importantbox}{本期核心命题}
具身智能不是“算法公司加一个机器人壳”，而是一条从整机、供应链、数据体系、AI infra、算法模型、分销渠道到客户价值的长链条。短期算法可快速传播，长期壁垒更可能来自真实世界闭环。
\end{importantbox}

\begin{knowledgebox}{视觉策略说明}
本视频是固定访谈画面，没有 slides、白板或产品演示。按播客笔记标准，正文不重复插入人物帧；封面用于来源识别，正文使用自动驾驶商业模式、数据飞轮、机器人大脑双系统等概念图来解释访谈内容。
\end{knowledgebox}

\subsection{本章小结}

本期不是简单的创始人故事，而是自动驾驶经验如何迁移到具身智能的案例。后文会依次讨论高继扬的方法论、Waymo 与 Momenta 的对照、星海图的整机供应链选择、Data Recipe、VLM/VLA 机器人大脑，以及“许华哲离开”背后的组织取舍。

